\section{Problem formulation and claim boundaries}

\subsection{Objects, metadata anchors, and the target question}

Let $s=(\text{state},\text{county},\text{site})$ denote a physical monitoring
site, let $p$ denote its reported POC, and let $y_{s,p,t}$ be the retained
daily concentration at date $t$. A monitor \emph{series} is the ordered pair
$(s,p)$; a physical site is $s$ alone. This distinction is essential because
different POCs can occur at the same location without furnishing independent
geographic information. For the primary parameter, an anchor is
\[
  a=(s,p,t_0,m^-,m^+),
\]
where $t_0$ is a persistent reported Method Code transition from $m^-$ to
$m^+$. The anchor is observable metadata, not a verified intervention.

For each anchor, the computational question is whether a pre-specified
cross-site residual changes around $t_0$ more than would be expected from the
benchmark's comparison procedures. It is \emph{not} ``did a device change?''
or ``what was the true pollution concentration?'' These questions would
require station histories, calibration records, deployment information, and
possibly meteorological and emissions evidence that do not follow from an AQS
Method Code field alone.

\subsection{Time-respecting residual estimand}

For a target series $i=(s,p)$ and qualified donor series $j\in N_a$, the
analysis uses the fixed transformation
\[
  z_{i,t}=\log(1+\max(y_{i,t},0)).
\]
With nonnegative donor weights $w_{ij}$ that sum to one, define the
calibration-centered residual
\[
 r_{i,t}=z_{i,t}-\sum_{j\in N_a}w_{ij}z_{j,t}
 -\operatorname{median}_{u\in T_{\mathrm{cal}}}
 \left(z_{i,u}-\sum_{j\in N_a}w_{ij}z_{j,u}\right).
\]
The protocol uses a 180-day calibration interval ending 15 days before the
anchor, a 60-day pre-anchor comparison interval, and a 60-day post-anchor
interval. It requires at least 30 retained observations in each required
window and at least two available donors on retained dates. The displayed
protocol-defined effect is
\[
  \widehat{\Delta}_a =
  \operatorname{median}_{t\in T_{\mathrm{post}}}r_{i,t}
  -\operatorname{median}_{t\in T_{\mathrm{pre}}}r_{i,t}.
\]
It is a local residual discontinuity on the log scale. It is not an estimate
of a causal physical bias or a corrected concentration.

\subsection{Claim boundaries}

\begin{table}[!htbp]
\centering
\caption{Claim boundaries imposed by the MetaShift-Bench protocol.}
\label{tab:claim-boundaries}
\scriptsize
\begin{tabular}{p{0.30\linewidth}p{0.30\linewidth}p{0.30\linewidth}}
\toprule
The evidence can support & The evidence does not establish & Required interpretation \\
\midrule
A reproducible metadata-anchored local residual audit & That a Method Code transition is a physical replacement or fault & Treat the transition as an anchor, not physical ground truth. \\
Benchmark-specific synthetic trade-offs among frozen methods & A robust aggregate MetaShift superiority claim & Report the paired intervals and retain the negative result. \\
Protocol-defined evidence tiers and abstentions & Instrument bias, true pollution correction, or causal attribution & Use tiers only to prioritize further records-based review. \\
Conditional interval behavior on frozen synthetic data & Coverage-calibrated confidence intervals for real-event physical bias & Describe real-event intervals as diagnostic. \\
\bottomrule
\end{tabular}
\par\smallskip
\begin{minipage}{0.94\linewidth}
\footnotesize\textit{Note.} These boundaries are design constraints, not post hoc caveats.
\end{minipage}
\end{table}


The benchmark's core logic is comparable to weakly labelled audit problems:
an observed administrative marker gives a reproducible condition to inspect,
but not a ground-truth class \citep{liu2016,frenay2014}. Accordingly, the
protocol makes three commitments. First, it fixes the anchor and the
time-respecting donor construction before the local residual is measured.
Second, it uses synthetic perturbations in stable regimes whenever an outcome
requires known truth. Third, it preserves non-complete cases as abstentions
rather than treating missing evidence as a negative result.

\subsection{What a negative model result means}

MetaShift is evaluated as a metadata-informed weighting hypothesis alongside
standard synthetic control, nearest-neighbor difference-in-differences, and
single-series change-point baselines. The benchmark permits a benchmark-specific
performance comparison, including paired uncertainty for local-effect error.
It does not permit a conclusion of robust superiority from a favorable point
estimate whose paired uncertainty interval crosses zero. This boundary protects
against selection after inspecting the held-out partition and keeps the paper's
main contribution focused on auditable benchmark construction.
